\documentclass[a4paper, 11pt]{article}
\usepackage{verbatim}
\usepackage[colorlinks]{hyperref}
\usepackage{graphicx, float}
\graphicspath{{assets/images/}}
% \usepackage{amsmath, amssymb, amsthm, amsfonts}
\usepackage{amsmath}
\usepackage{tcolorbox}
\tcbset{colback=blue!4!white,colframe=blue!50!green!30!gray}
% Margins
\usepackage[top=1.5cm, bottom=2cm, left=2.4cm, right=1.8cm]{geometry}
% Temporary fix until better understanding of the indentation rules
\setlength{\parindent}{0pt}

\title{Yet Another \LaTeX{} Tutorial \\
  \medskip \large A road to Procrastination}
\author{Energos}
\date{March 2025}

\begin{document}

\maketitle

\section{Why?}
\href{https://www.youtube.com/watch?v=9eLjt5Lrocw}
{Why \LaTeX{}? - YouTube}

\section{Install \LaTeX{} in Debian}
\begin{verbatim}
apt install texlive texlive-science texlive-lang-english texlive-xetex
\end{verbatim}

\section{The Document}
\begin{verbatim}
\begin{document}
\maketitle
\section{Section title}
\subsection{Subsection title}
\end{document}
\end{verbatim}

\section{Title}
\begin{verbatim}
\title{Yet Another \LaTeX{} Tutorial \\
       \medskip \large A road to Procrastination}
\author{Energos}
\date{March 2025}
\end{verbatim}

\section{A Section}
\begin{verbatim}
\section{A Section}
\end{verbatim}

\subsection{A Subsection}
\begin{verbatim}
\section{A Subsection}
\end{verbatim}

\section{Lorem ipsum}

Lorem ipsum dolor sit amet. Qui officiis repellendus id sunt consequatur
ut nisi eligendi est rerum dolorem qui rerum porro non laudantium quasi ut
laudantium aspernatur.

Sed nesciunt voluptatem eos internos delectus sed blanditiis explicabo nam
officia voluptatem vel provident minus non distinctio nihil.

\section{Paragraph indentation}
Temporary fix until better understanding of the indentation rules.
\begin{verbatim}
\setlength{\parindent}{0pt}
\end{verbatim}

\section{Margins}
\begin{verbatim}
\usepackage[top=0.75in, bottom=1.00in, left=1.20in, right=1.00in]{geometry}
\end{verbatim}

\section{Lists}
\begin{itemize}
\item One
\item Two
\item Three
\end{itemize}

\begin{verbatim}
\begin{itemize}
\item One
\item Two
\item Three
\end{itemize}
\end{verbatim}

\section{URL links}
\href{https://www.ctan.org/pkg/hyperref}{CTAN: Package hyperref}

\begin{verbatim}
\usepackage[colorlinks]{hyperref}
\href{https://www.ctan.org/pkg/hyperref}{CTAN: Package hyperref}
\end{verbatim}

\section{Verbatim}
\begin{verbatim}
\usepackage{verbatim}
\end{verbatim}

\subsection{Including verbatim text}

\verb|\begin{verbatim}| \\
\verb|Verbatim stuff.| \\
\verb|More verbatim stuff...| \\
\verb|\end{verbatim}|

\subsection{Including a file verbatim}
\begin{verbatim}
\verbatiminput{tutorial000.tex}
\end{verbatim}

\section{Images}
\begin{figure}[H]
  \centering
  \includegraphics[width=4in]{latex-project.png}
  \caption{\LaTeX{}}
  \label{fig:latex-project-logo}
\end{figure}

\begin{verbatim}
\usepackage{graphicx, float}
\graphicspath{{assets/images/}}

\begin{figure}[H]
  \centering
  \includegraphics[width=4in]{latex-project.png}
  \caption{\LaTeX{}}
  \label{fig:latex-project-logo}
\end{figure}
\end{verbatim}

SVG Images require more work:
\href{https://www.baeldung.com/cs/latex-svg-images}
{How to Use SVG Images in \LaTeX{}}.

\section{Equations}

\begin{align}
  \nabla \cdot \mathbf{E} \,\,\, &= \frac{\rho}{\varepsilon_0} \\
  \nabla \cdot \mathbf{B} \,\,\, &= 0 \\
  \nabla \times \mathbf{E} &= -\frac{\partial \mathbf{B}}{\partial t} \\
  \nabla \times \mathbf{B} &= \mu_0 \left(\mathbf{J}
                             + \varepsilon_0 \frac{\partial
                             \mathbf{E}}{\partial t} \right)
\end{align}

\begin{align}
  \sin x &= \sum^{\infty}_{n=0} \frac{(-1)^n}{(2n+1)!} x^{2n+1}
  &= x - \frac{x^3}{3!} + \frac{x^5}{5!} - \frac{x^7}{7!} + \frac{x^9}{9!}
    - \cdots && \text{for all } x\\[6pt]
  \cos x &= \sum^{\infty}_{n=0} \frac{(-1)^n}{(2n)!} x^{2n}
  &= 1 - \frac{x^2}{2!} + \frac{x^4}{4!} - \frac{x^6}{6!} + \frac{x^8}{8!}
    - \cdots && \text{for all } x
\end{align}

\begin{align}
  A &= 4 \pi r^2 \\
  V &= \frac{4}{3} \pi r^3
\end{align}

\begin{verbatim}
\usepackage{amsmath}

\begin{align}
  A &= 4 \pi r^2 \\
  V &= \frac{4}{3} \pi r^3
\end{align}
\end{verbatim}

\section{Simple Tab Stops}
\begin{tabbing}
  ~~~~~~~~~~~~~~ \= ~~~~~~~~~~~~~~ \= 0 \\
  1 \> 1 \> 1 \\
  22 \> 22 \> 22 \\
  333 \> 333 \> 333 \\
  4444 \> 4444 \> 4444
\end{tabbing}

\begin{verbatim}
\begin{tabbing}
  ~~~~~~~~~~~~~~ \= ~~~~~~~~~~~~~~ \= 0 \\
  1 \> 1 \> 1 \\
  22 \> 22 \> 22 \\
  333 \> 333 \> 333 \\
  4444 \> 4444 \> 4444
\end{tabbing}
\end{verbatim}

\section{Boxes}
Draw colored and framed boxes using
\href{https://ctan.org/pkg/tcolorbox?lang=en}{tcolorbox}.
This package has a huge set of options and features.
Please read the
\href{http://mirrors.ctan.org/macros/latex/contrib/tcolorbox/tcolorbox.pdf}
{documentation} and
\href{http://mirrors.ctan.org/macros/latex/contrib/tcolorbox/tcolorbox-example.pdf}
{examples of use}.

\begin{tcolorbox}
  This is a box using document wide default colors.
\end{tcolorbox}

\begin{tcolorbox}[colback=red!30!white,colframe=red!80!black]
  This is a somewhat reddish box.
\end{tcolorbox}

\begin{tcolorbox}
\begin{verbatim}
\usepackage{tcolorbox}
% Define document wide default colors
\tcbset{colback=blue!4!white,colframe=blue!50!green!30!gray}

\begin{tcolorbox}
  This is a box using document wide default colors.
\end{tcolorbox}

\begin{tcolorbox}[colback=red!30!white,colframe=red!80!black]
  This is a somewhat reddish box.
\end{tcolorbox}
\end{verbatim}
\end{tcolorbox}

\section{Packages}
- \href{https://www.ctan.org/pkg/hyperref}{CTAN: Package hyperref} \\
- \href{https://ctan.org/pkg/verbatim}{CTAN: Package verbatim} \\
- \href{https://ctan.org/pkg/geometry}{CTAN: Package geometry} \\
- \href{https://www.ctan.org/pkg/graphicx}{CTAN: Package graphicx} \\
- \href{https://www.ctan.org/pkg/float}{CTAN: Package float} \\
- \href{https://ctan.org/pkg/amsmath?lang=en}{CTAN: Package amsmath} \\
- \href{https://ctan.org/pkg/tcolorbox?lang=en}{CTAN: Package tcolorbox}

\section{Generate a .PDF file}
\begin{verbatim}
xelatex tutorial000.tex
\end{verbatim}
or
\begin{verbatim}
pdflatex tutorial000.tex
\end{verbatim}

\section{Resources}
- \href{https://www.latex-project.org/}
{LaTeX - A document preparation system} \\
- \href{https://www.ctan.org/}{CTAN: Comprehensive TeX Archive Network} \\
- \href{https://texdoc.org/index.html}
{texdoc - TeX and LaTeX documentation lookup system} \\
- \href{https://www.ams.org/arc/resources/amslatex-about.html}
{American Mathematical Society - LaTeX extensions} \\
- \href{https://www.ams.org/arc/resources/texresources.html}
{American Mathematical Society - TeX Resources} \\
- \href{https://texfaq.org/}
{The TeX Frequently Asked Question List - The TeX FAQ} \\
- \href{https://texfaq.org/FAQ-latex-books}{Books on LaTeX - The TeX FAQ} \\
- \href{https://www.gnu.org/software/auctex/}
{AUCTeX - Sophisticated document creation in GNU Emacs}

\section{Tutorials}
- \href{https://www.youtube.com/watch?v=lgiCpA4zzGU}
{LaTeX for Students – A Simple Quickstart Guide - YouTube} \\
- \href{https://www.youtube.com/watch?v=VqMETR1GeYI}
{LaTeX Basics 00: Installation and Hello, World! - YouTube} \\
- \href{https://www.youtube.com/watch?v=_q4aJdbKXKI}
{LaTeX Basics 0½: Emacs - YouTube}

\section{The whole enchilada}
\verbatiminput{tutorial000.tex}

\end{document}
